
\only<presentation>{
\againframe{History}
\againframe{CourseStructure}
\againframe{CloneTheCourse}

\againframe{CourseOverview}
\againframe{RecapDescriptiveStatistics}
\againframe{RecapProbabilityFundamentals}
\againframe{RecapDiscreteDistributions}
\againframe{RecapContinousProbabilityDistributions}

\begin{frame}[allowframebreaks]{Recap: Linear Regression}
    We moved from analysing single variables in isolation to measuring the interaction and predictive relationships between two variables.

    \textbf{1. Correlation ($r$)}
    \begin{itemize}
        \item \textbf{Strength \& Direction:} Measures linear connection between $-1$ and $+1$.
        \item \textbf{The Golden Rule:} Correlation does \textbf{not} imply causation; beware of \lq lurking variables\rq\ and the Storks Paradox.
    \end{itemize}

    \textbf{2. Simple Linear Regression}
    \begin{itemize}
        \item \textbf{Predictive Model:} $\hat{y} = mx + c$, where $m$ is the system \textbf{Sensitivity} and $c$ is the \textbf{Zero-Offset} error.
        \item \textbf{Least-Squares Method:} The math calculates the line that minimises the sum of the \textbf{squared residuals} (vertical errors).
    \end{itemize}

    \textbf{3. Model Evaluation \& Mapping}
    \begin{itemize}
        \item \textbf{Goodness of Fit ($R^2$):} The percentage of the output's variation that is explicitly explained by our mathematical model.
        \item \textbf{Non-Linear Mapping:} Using natural logs ($\ln$) to transform Exponential and Power laws into linear space for easy calculation.
    \end{itemize}

    \begin{alertblock}{Course Conclusion}
        You now have the tools to describe past data, predict future probabilities, and model complex engineering interactions. Good luck with your revision!
    \end{alertblock}
\end{frame}

\againframe{Prize}

\againframe{StudentAcknowledgements}

\againframe{Assessment}
\begin{frame}{Examination Strategy and Standards}
    \begin{itemize}
        \item \textbf{Structure:} The examination comprises \textbf{4 independent questions} in STACK.
        \item \textbf{Mandatory Precision:} All numerical results \textit{must} be reported to \textbf{4 significant figures (s.f.)}.
        \item \textbf{Curriculum Alignment:} 3 questions are directly modelled on the \textbf{Section B Exercise Sheets}.
        \item \textbf{Abstract Application:} 1 question involves a conceptual \lq twist\rq\ .
        \begin{itemize}
            \item Specifically testing your ability to interpret parameters in a novel context.
            \item Focus is on \textbf{Engineering Intuition}.
        \end{itemize}
    \end{itemize}
\end{frame}
}